% BHU PYQ -- Auto-generated & error-corrected
\documentclass[11pt,a4paper]{article}

\usepackage[a4paper,top=2.5cm,bottom=2.5cm,left=2.8cm,right=2.8cm,headheight=14pt]{geometry}
\usepackage{amsmath,amssymb}
\usepackage[T1]{fontenc}
\usepackage[utf8]{inputenc}
\usepackage{lmodern}
\usepackage{microtype}
\usepackage{fancyhdr}
\usepackage{enumitem}
\usepackage{hyperref}
\usepackage{lastpage}
\usepackage{parskip}

\hypersetup{colorlinks=true,urlcolor=black,linkcolor=black,
  pdftitle={STB-301: Probability Theory and Statistical Methods -- BHU PYQ},pdfauthor={Banaras Hindu University}}

\pagestyle{fancy}\fancyhf{}
\renewcommand{\headrulewidth}{0.4pt}\renewcommand{\footrulewidth}{0.4pt}
\fancyhead[L]{\small\textit{Banaras Hindu University}}
\fancyhead[R]{\small\textit{STB-301: Probability Theory and Statistical }}
\fancyfoot[C]{\small Page \thepage\ of \pageref{LastPage}}

\newlist{parts}{enumerate}{2}
\setlist[parts,1]{label=(\alph*),leftmargin=*,itemsep=5pt,topsep=3pt}
\setlist[parts,2]{label=(\roman*),leftmargin=*,itemsep=3pt,topsep=2pt}
\newcommand{\pts}[1]{\hfill\textit{[\#1]}}

\begin{document}

\begin{center}
  {\large\bfseries BANARAS HINDU UNIVERSITY}\\[2pt]
  {\normalsize B.Sc. (Hons.) Semester III Examination, 2022-23}\\[6pt]
  \rule{0.8\linewidth}{0.5pt}\\[6pt]
  {\large\bfseries Paper No. STB-301: Probability Theory and Statistical Methods}\\[4pt]
  {\normalsize Time: 3 Hours\hfill Maximum Marks: 70}
\end{center}
\rule{\linewidth}{0.4pt}
\smallskip
\noindent\textit{Instructions: Answer \textbf{five} questions including Question~1 which is compulsory. Symbols have their usual meanings.}
\bigskip

\medskip\noindent\textbf{Question 1.}
\begin{parts}
  \item What do you mean by a random sample from a distribution? Distinguish between standard deviation and standard error. Show that the sampling distribution of the sum of two independent Poisson variates is again a Poisson variate. \pts{7}
  \item What are the criteria of a good estimator? Explain any two of them with examples. Let $X_1, X_2, \dots, X_n$ be a random sample from a $N(\mu, \sigma^2)$ population, and $T = \sum a_i X_i$ be the statistic based on sample values. Show that:
  \begin{parts}
    \item $E(T) = \mu$ when $\sum a_i = 1$.
    \item $V(T) = \sigma^2 \sum a_i^2$.
  \end{parts} \pts{7}
\end{parts}

\medskip\rule{\linewidth}{0.2pt}

\medskip\noindent\textbf{Question 2.} Define null and alternative hypotheses with a suitable example. Explain the concept of Type-I and Type-II errors and bring out their importance in a testing of hypothesis problem. A random sample of size 1 is taken from a Poisson distribution with parameter $\lambda$. If $x \ge 2$ is the critical region for testing the hypothesis $H_0: \lambda = 2$ against $H_1: \lambda = 1$, find the size and power of the test. \pts{14}

\medskip\rule{\linewidth}{0.2pt}

\medskip\noindent\textbf{Question 3.}
\begin{parts}
  \item Define $\chi^2$ statistic. Describe the $\chi^2$ test for testing independence of attributes. \pts{7}
  \item What is the $F$-distribution? Explain the testing procedure for testing the equality of two variances from two independent normal populations. \pts{7}
\end{parts}

\medskip\rule{\linewidth}{0.2pt}

\medskip\noindent\textbf{Question 4.} State the test of significance for the difference between two means from two independent normal populations when the variances are equal but unknown. Below are given the gains in weights (in kgs) of pigs fed on two diets $A$ and $B$:
\begin{center}
\begin{tabular}{ll}
Diet $A$: & 32, 35, 33, 29, 24, 41, 42, 36 \\
Diet $B$: & 35, 36, 39, 37, 38, 25, 26, 30, 24, 20
\end{tabular}
\end{center}
Test if the two diets differ significantly as regards their effect on increase in weight at $5\%$ level of significance, given that $t_{0.025, 16} = 2.12$. \pts{14}

\medskip\rule{\linewidth}{0.2pt}

\medskip\noindent\textbf{Question 5.}
\begin{parts}
  \item Describe the importance of the Central Limit Theorem. State the test of significance for the difference between two proportions in the case of a large sample. \pts{7}
  \item Describe the $\chi^2$ goodness-of-fit test, stating its underlying assumptions. \pts{7}
\end{parts}

\medskip\rule{\linewidth}{0.2pt}

\medskip\noindent\textbf{Question 6.} Define order statistics along with their various applications. Show that the distribution of the $r$-th order statistic from a uniform distribution $U(0, 1)$ follows a beta distribution. Also, deduce the distribution of the smallest and largest order statistics. \pts{14}

\medskip\rule{\linewidth}{0.2pt}

\medskip\noindent\textbf{Question 7.}
\begin{parts}
  \item Explain the advantages and disadvantages of non-parametric tests. Describe the one-sample Kolmogorov-Smirnov non-parametric test. \pts{7}
  \item Develop the median non-parametric test, stating the underlying assumptions and the null hypothesis. \pts{7}
\end{parts}

\medskip\rule{\linewidth}{0.2pt}

\medskip\noindent\textbf{Question 8.} Write short notes on any three of the following:
\begin{enumerate}[label=\roman*.]
  \item Estimate and estimator
  \item Sampling distribution
  \item Simple and composite hypothesis with examples
  \item Fisher's $t$-distribution and its uses
  \item Fisher's $z$-transformation
\end{enumerate} \pts{14}

\vfill
\rule{\linewidth}{0.4pt}
\begin{center}\small
  \href{https://bhu-exam.vercel.app/}{Click here to visit our website}\\[4pt]\href{https://me-aryan.vercel.app/}{Click here to visit our contributor}\\[4pt]\href{https://quashan-me.vercel.app/}{Click here to visit our contributor}
\end{center}

\end{document}